\selectlanguage{ngerman}
\begin{ledgroupsized}[r]{120mm}
\footnotesize
\pstart
\noindent\textbf{Überlieferung:}
\pend
\end{ledgroupsized}
\begin{ledgroupsized}[r]{114mm}
\footnotesize
\pstart \parindent -6mm
\makebox[6mm][l]{\textit{L}}%
Aufzeichnung:
LH~XXXV~9,\;7 Bl.\;22\textendash23. 
Ein Bogen~8\textsuperscript{o},
Ränder beschnitten;
Fragment eines \href{\Scan{LH_35_09_07_022-023wz}}{Wasserzeichens}: im Falz:
Niedersächsiches Papier.
Zweieinviertel Seiten;
Bl.\;23\,v\textsuperscript{o} ist unbeschrieben.
\pend
\end{ledgroupsized}
\vspace{3.5mm}
\begin{ledgroup}
\footnotesize
\pstart
\noindent%
\textbf{Datierungsgründe:}
Die Datierung der Aufzeichnung N.\;\ref{41121} beruht auf äußeren Eigenschaften des Text\-trägers %
sowie auf inhaltlichen Betrachtungen.
Siehe die ausführliche Begründung auf S.\;\pageref{41121_dat}.
\pend
\end{ledgroup}
\selectlanguage{latin}
\frenchspacing%
\vspace{5mm}
\pstart%
\normalsize%
\noindent%
\lbrack%
\href{\Scan{LH_35_09_07_022r}}%
{22\,r\textsuperscript{o}}%
\rbrack\ %
Si duo corpora%
\indexS{corpora occurrentia sibi directe centraliter}
(quae instar punctorum considerabimus)%
\indexS{corpus instar puncti}
sibi occurrant directe%
\indexS{occursus directus centralis}
\edtext{centraliter,
has}{%
\lemma{centraliter,}%
\Bfootnote{%
\textit{(1)}~et
\textit{(2)}~has%
~\textit{L}}}
habemus regulas.%
\indexS{regulae de concursu corporum}
(1)
Eadem est celeritas ad se%
\indexS{celeritas accedendi ante concursum}
\edtext{invicem accedendi}{%
\lemma{invicem}\Bfootnote{%
\textit{(1)}~recedendi
\textit{(2)}~accedendi%
~\textit{L}}}
ante concursum,
et a se invicem recedendi%
\indexS{celeritas recedendi post concursum}
post concursum.%
\indexS{concursus directus centralis}
(2)
Eadem est quantitas progressus%
\indexS{quantitas progressus ante concursum}%
\indexS{quantitas progressus post concursum}
ante et post
\edtext{%
concursum in easdem partes,
vel quod eodem redit,
eadem est celeritas centri gravitatis.%
\indexS{centrum gravitatis commune}%
\indexS{celeritas centri gravitatis}
(3)
Eadem%
}{%
\lemma{concursum}\Bfootnote{%
\textit{(1)}~. (3) Ea %
\textit{(2)}~in easdem partes %
\textit{(a)}~. (3) Eadem %
\textit{(b)}~, vel quod \lbrack...\rbrack\ (3) Eadem%
~\textit{L}}}
est quantitas virium.%
\indexS{quantitas virium in concursu}
(4)
Si corpora tendant ad easdem partes ante et post concursum,%
\indexS{corpora tendentia in easdem partes}
\edtext{%
servatur quantitas motus%
\indexS{quantitas motus servata}
seu summa motuum;%
\indexS{summa motuum}
si%
}{%
\lemma{servatur}\Bfootnote{%
\textit{(1)}~eadem quantitas mo
\textit{(2)}~summa
\textit{(a)}~motuum; si
\textit{(b)}~quantitatum motus
\textit{(3)}~quantitas motus %
\lbrack...\rbrack\ motuum; si%
~\textit{L}}}
corpora tendant in diversas partes ante et post concursum,%
\indexS{corpora tendentia in diversas partes}
servatur differentia%
\indexS{differentia motuum}
\edtext{%
motuum;
denique si corpora%
}{%
\lemma{motuum}\Bfootnote{%
\textit{(1)}~. (5) Si co %
\textit{(2)}~; denique si corpora%
~\textit{L}}}
tendant in partes easdem ante ictum,%
\indexS{corpora tendentia in easdem partes}%
\indexS{ictus}
in contrarias post ictum%
\indexS{corpora tendentia in contrarias partes}%
\indexS{ictus}%
\lbrack,\rbrack\
quae erat summa motuum%
\indexS{summa motuum}
fit differentia;%
\indexS{differentia motuum}
si tendant in contrarias ante ictum,%
\indexS{corpora tendentia in contrarias partes}%
\indexS{ictus}
in easdem post ictum,%
\indexS{corpora tendentia in easdem partes}%
\indexS{ictus}
ex differentia motuum%
\indexS{differentia motuum}
fit summa.%
\indexS{summa motuum}
\pend%
\vspace{0.25em}%
\pstart%
Calculus%
\indexS{calculus}
\pend%
\vspace{0.50em}%
\pstart%
\noindent%
Sit corpus
$\displaystyle\efrac{a}{b},$
celeritas ante ictum%
\indexS{celeritas ante ictum}%
\indexS{ictus}
$\displaystyle\efrac{v}{y},$
post ictum%
\indexS{celeritas post ictum}%
\indexS{ictus}
$\displaystyle\efrac{x}{z}.$%
\edlabel{35_09_07_022-023_1a}%
\edtext{}{%
\xxref{35_09_07_022-023_1a}{35_09_07_022-023_1b}%
\lemma{$\displaystyle\efrac{x}{z}.$\rule[-2ex]{0pt}{5ex}}%
\Bfootnote{%
\lbrack/\rbrack\
\textit{(1)}~$v \PM y$
\textit{(2)}~%
Sit%
~\textit{L}%
}}%
\pend%
\vspace{0.75em}%
\pstart%
\noindent%
\lbrack\textit{Nachfolgend kleingedruckter Text in L gestrichen:}\rbrack%
\pend%
\vspace{0.5em}%
\footnotesize%
\pstart%
\noindent%
Sit autem \textit{a}%
\edlabel{35_09_07_022-023_1b}
\edtext{%
corpus,
quod directioni contrariae non succumbit,%
\indexS{corpus succumbens directioni contrariae}%
\indexS{directio contraria}
seu cujus quantitas motus non est minor,%
\indexS{quantitas motus}%
}{%
\lemma{corpus,}%
\Bfootnote{%
\textit{(1)}~cujus quantitas motus est major, quodque adeo contrariae non succumbit,%
\indexS{corpus succumbens quantitati motus contrariae}
\textit{(2)}~quod directioni %
\lbrack...\rbrack\ cujus quantitas
\textit{(a)}~motus \textlangle prae\textrangle\
\textit{(b)}~motus non est minor,%
~\textit{L}%
}}
seu quod in easdem partes tendit cum centro%
\indexS{corpus tendens in easdem partes centri gravitatis}%
\indexS{centrum gravitatis commune}
\edtext{%
gravitatis.%
\quad%
$v - y \protect\overset{(1)}{=} \PM\, x \MP z,$%
\rule[-2ex]{0pt}{5ex}
explicando literam%
}{%
\lemma{gravitatis.}%
\Bfootnote{\lbrack/\rbrack\
\textit{(1)}~$v \PM y \protect\overset{(1)}{=} x - z$
\textit{(2)}~$v - y \protect\overset{(1)}{=} \PM\, x \MP z,$\rule[-2ex]{0pt}{5ex}
\textit{(a)}~pre
\textit{(b)}~sumendo lit
\textit{(c)}~explicando literam%
~\textit{L}%
}}
per numerum negativum,%
\indexS{litera explicanda per numerum negativum}%
\indexS{numerus negativus}
cum motus \textit{y} est contrarius ipsi \textit{v}.%
\indexS{motus contrarius}
\pend%
\newpage%
\normalsize%
\pstart%
\noindent%
Sit autem
$\displaystyle\efrac{a}{b}$
corpus quod%
\indexS{corpus antecedens centrum gravitatis}%
\indexS{corpus sequens centrum gravitatis}
$\displaystyle\efrac{\text{sequitur}}{\text{antecedit}}$
centrum gravitatis%
\indexS{centrum gravitatis commune}
ante ictum.%
\indexS{ictus}%
\pend%
\vspace{0.75em}%
\pstart%
\noindent%
Quodsi ambo tendant in easdem partes,
simul et \textit{a} sequetur et \textit{b} antecedet%
\lbrack;\rbrack\
quodsi tendant in contrarias partes,
tamen
\edtext{%
semper \textit{a} esto%
}{%
\lemma{semper \textit{a}}%
\Bfootnote{%
\textit{(1)}~erit
\textit{(2)}~esto%
~\textit{L}%
}}
illud
quod sequitur centrum
\edtext{%
gravitatis,%
\indexS{centrum gravitatis commune}
et \textit{b} centrum antecedet.
Potest tamen%
}{%
\lemma{gravitatis,}%
\Bfootnote{%
\textit{(1)}~et tunc e
\textit{(2)}~et \textit{b} %
\lbrack...\rbrack\ Potest tamen%
~\textit{L}%
}}
fieri
ut quod in easdem cum centro partes tendit
reflectatur.%
\indexS{corpus reflexum post ictum}
In
\edtext{%
casu motus contrarii%
\indexS{motus contrarius}%
}{%
\lemma{casu}%
\Bfootnote{%
\textit{(1)}~occursus
\textit{(2)}~motus contrarii%
~\textit{L}%
}}
corpus cujus
\edlabel{35_09_07_022-023_4a}%
quantitas motus%
\indexS{quantitas motus}%
\edtext{}{%
\xxref{35_09_07_022-023_4a}{35_09_07_022-023_4b}%
\lemma{quantitas motus}%
\Bfootnote{%
\bars{seu \textit{streicht Hrsg.}}
\textit{(1)}~\href{\Scan{LH_35_09_07_022v}}{\lbrack22~v\textsuperscript{o}\rbrack} seu quantitas progressus major est, in easdem partes tendit cum centro gravitatis. Sit \textit{a}
\textit{(2)}~\href{\Scan{LH_35_09_07_022r}}{\lbrack22~r\textsuperscript{o}\rbrack} seu quantitas \lbrack...\rbrack\ easdem partes
\textit{(a)}~centro gravitatis
\textit{(b)}~tendit cum centro gravitatis. \href{\Scan{LH_35_09_07_022v}}{\lbrack22~v\textsuperscript{o}\rbrack} \lbrack/\rbrack\ Generaliter%
~\textit{L}%
}}
seu quantitas progressus%
\indexS{quantitas progressus}
major est%
\lbrack,\rbrack\
in easdem partes tendit cum centro gravitatis.%
\indexS{corpus tendens in easdem partes centri gravitatis}%
\indexS{centrum gravitatis commune}
\lbrack%
\href{\Scan{LH_35_09_07_022v}}%
{22\,v\textsuperscript{o}}%
\rbrack%
\pend%
\vspace{0.5em}%
\pstart%
\edlabel{LH_35_09_07_022-023_generaliter-1}%
Generaliter%
\edlabel{35_09_07_022-023_4b}%
\lbrack:\rbrack\
\pend%
\vspace{0.25em}%
\pstart%
\noindent%
\protect%
\hspace*{30.5mm}$avv + byy$%
\hspace{3.5mm}$\overset{(1)}{=}$%
\hspace{4.8mm}$axx + bzz$%
\hspace{4mm}per reg. 3%
\newline%
\protect%
\hspace*{32.0mm}$av + by$%
\hspace{4.5mm}$\overset{(2)}{=}$%
\hspace{6.9mm}$ax + bz$%
\hspace{5.6mm}per reg. 2%
\newline%
\edlabel{35_09_07_022-023_3a}%
\protect%
ex 1 fit%
\hspace{19.5mm}$a\,\overline{vv - xx}$%
\hspace{4mm}$\overset{(3)}{=}$%
\hspace{5.1mm}$b\,\overline{zz - yy}$%
\newline%
\edtext{}{%
\xxref{35_09_07_022-023_3a}{35_09_07_022-023_3b}%
\lemma{\textit{Am Rand, %
gestrichen:}}%
\Afootnote{%
{\footnotesize%
Corpus
quod ante ictum%
\indexS{ictus}
in easdem partes tendit cum centro grav.%
\indexS{corpus tendens in easdem partes centri gravitatis}
vocamus \textit{a},
alterum \textit{b}.}%
}}%
\protect%
ex 2 fit%
\hspace{21.5mm}$a\,\overline{v - x}$%
\hspace{6.2mm}$\overset{(4)}{=}$%
\hspace{6.9mm}$b\,\overline{z - y}$%
\newline%
\protect%
divid. 3 per 4, fit%
\hspace{8.0mm}$v + x$%
\hspace{6.2mm}$\overset{(5)}{=}$%
\hspace{8.5mm}$y + z$%
\newline%
\protect%
ex 5%
\hspace{28.5mm}$v - y$%
\hspace{6.3mm}$\overset{(6)}{=}$%
\hspace{8.2mm}$z-x.$%
\edlabel{35_09_07_022-023_3b}%
\hspace{6.0mm}Est
\edlabel{35_09_07_022-023_2a}%
Reg.\;1.%
\edtext{}{%
\xxref{35_09_07_022-023_2a}{35_09_07_022-023_2b}%
\lemma{Reg.\;1.}%
\Bfootnote{%
\lbrack/\rbrack\
\textit{(1)}~Corporis $\displaystyle\efrac{a}{b}$
\textit{(a)}~celeritas
\textit{(b)}~progressus ante ictum $\displaystyle\efrac{v}{y}$ post ictum $\displaystyle\efrac{x}{z}.$
\textit{(2)}~Progressus cum \lbrack...\rbrack\ ictum $\displaystyle\efrac{x}{z}.$%
~\textit{L}%
}}%
\pend%
\vspace{0.75em}%
\pstart%
\noindent%
Progressus cum centro gravitatis%
\indexS{progressus cum centro gravitatis communi}
corporis $\displaystyle\efrac{a}{b}$
ante ictum $\displaystyle\efrac{v}{y}$
post ictum $\displaystyle\efrac{x}{z}.$%
\edlabel{35_09_07_022-023_2b}%
\edlabel{LH_35_09_07_022-023_generaliter-2}%
\pend%
\vspace{0.5em}%
\pstart%
\noindent%
\edtext{%
Quodsi corpus%
}{%
\lemma{Quodsi}%
\Bfootnote{%
\textit{(1)}~quod
\textit{(2)}~corpus%
~\textit{L}%
}}
moveatur in partes centro gravitatis,%
\indexS{centrum gravitatis commune}
vel praevalenti directioni%
\indexS{directio praevalens in concursu}%
\lbrack,\rbrack\
contrarias,%
\indexS{corpus movens contra centrum gravitatis}%
\indexS{corpus movens contra directionem praevalentem}
ejus litera explicanda est%
\indexS{litera explicanda per numerum negativum}
per numerum quidem graduum celeritatis,%
\indexS{numerus graduum celeritatis}%
\indexS{gradus celeritatis}%
\indexS{celeritas ante ictum}%
\indexS{celeritas post ictum}
sed praefixo
\edlabel{35_09_07_022-023_5a}%
\edtext{}{%
\xxref{35_09_07_022-023_5a}{35_09_07_022-023_5b}%
\lemma{signo $-$.}%
\Bfootnote{%
\textit{(1)}~Sit jam corpus \textit{a}, %
\textit{(2)}~Videamus jam \lbrack...\rbrack\ idem corpus%
~\textit{L}%
}}%
signo $-.$%
\indexS{signum praefixum}
\pend%
\pstart%
Videamus jam
an super idem corpus%
\edlabel{35_09_07_022-023_5b}
eodem modo ante concursum motum,
diversa idem efficere%
\indexS{effectus corporum diversorum super idem}
\edtext{%
possint.
Contenti autem simus aeqq. 4 et 5%
\indexS{aequatio}
cum caeterae inde sequantur;
et ponamus dari%
}{%
\lemma{possint.}%
\Bfootnote{%
\textit{(1)}~Ponantur
\textit{(2)}~Contenti autem %
\lbrack...\rbrack\ ponamus dari%
~\textit{L}%
}}
\textit{b}, \textit{y}, \textit{z},
quaeri \textit{a} et \textit{v},
eaque diversis modis.
Per aeq. 5 est
$v \overset{(7)}{=} y + z - x,$
et per 4 est
$v \overset{(8)}{=} \overline{b : a}\; \overline{z - y} + x.$
Per 7 et 8 fit
$\overline{b : a} \overset{(9)}{=} y + z - 2x,\, : \overline{z - y}.$
Ergo fiet:
$a\overset{(10)}{=} b\; \overline{z - y}\, :\, \overline{y - 2x + z}.$
Unde patet infinitis modis rem posse praestari,
prout variatur \textit{x}. %
Sed veniamus ad exemplum.%
\indexS{exemplum}
\pend%
\pstart
Sit
$b \overset{(11)}{=} 4,$
$z \overset{(12)}{=} 1,$
fiet:
$v \overset{(13)}{=} y + 1 - x,$
$a \overset{(14)}{=} 4 \cdot \overline{1 - y} : \overline{y - 2x + 1}.$
Si ponamus
$x \overset{(15)}{=} 0$
fiet
$v \overset{(16)}{=} y + 1$
et
$a \overset{(17)}{=} 4 \cdot \overline{1 - y}\, :\, \overline{ 1 + y}.$
Unde si sumatur
\edtext{%
$y \overset{(18)}{=} 1 : 2$
et
$x = 0,$
fiet%
}{%
\lemma{$y \protect\overset{(18)}{=} 1 : 2$}%
\Bfootnote{%
\textit{(1)}~fiet
\textit{(2)}~et $x = 0,$ fiet%
~\textit{L}%
}}
$v \overset{(19)}{=} 3 : 2$
et
$a \overset{(20)}{=} 4 : 3.$
Itaque corpus $4 : 3$
celeritate $3 : 2$
et corp. 4
celer. $1 : 2$
tendant in easdem%
\indexS{corpora tendentia in easdem partes}%
\indexS{celeritas ante ictum}
\edtext{%
partes,
post concursum%
}{%
\lemma{partes,}%
\Bfootnote{%
\textit{(1)}~concursu
\textit{(2)}~post concursum%
~\textit{L}%
}}
corpus $4 : 3$ quiescet%
\indexS{corpus quiescens post concursum}
et 4 perget celeritate 1.%
\indexS{corpus pergens post concursum}
Rursus retentis aeqq. 11. 12. 18%
\indexS{aequatio}
sit
$x \overset{(21)}{=} 1 : 2,$
fiet
$v \overset{(22)}{=} 1$
et
$a\overset{(23)}{=}4.$
Hoc est,
si corpus 4 celeritate 1,
et corpus 4 celeritate $1 : 2$
tendunt in easdem partes,
post concursum%
\indexS{corpora tendentia in easdem partes}%
\indexS{celeritas post ictum}
\edtext{%
corpus illud perget%
}{%
\lemma{corpus}%
\Bfootnote{%
\textit{(1)}~prius 4 p %
\textit{(2)}~illud perget%
~\textit{L}%
}}
celeritate $1 : 2,$
hoc celeritate ut 1.
\pend%
\pstart%
Ex quo exem%
\lbrack%
\href{\Scan{LH_35_09_07_023r}}{23\,r\textsuperscript{o}}%
\rbrack%
plo patet%
\indexS{exemplum}
\edtext{%
eleganter
super corpore \textit{A}, 4
moto celeritate $1 : 2$%
\indexS{celeritas ante ictum}
duo corpora%
}{%
\lemma{eleganter}%
\Bfootnote{%
\textit{(1)}~duo corpora
\textit{(2)}~super corpore %
\lbrack...\rbrack\ duo corpora%
~\textit{L}%
}}
ipsum persequentia,%
\indexS{corpus persequens}
\edtext{%
unum \textit{C}, $4 : 3$ motum%
}{%
\lemma{unum}%
\Bfootnote{%
\textit{(1)}~motum
\textit{(2)}~\textit{C}
\textit{(3)}~\textit{B}
\textit{(4)}~\textit{C}, $4 : 3$ motum%
~\textit{L}%
}}
celeritate $3 : 2,$
alterum \textit{B}, 4
motum celeritate 1,%
\indexS{celeritas ante ictum}
idem esse effectura,%
\indexS{effectus corporum diversorum super idem}
alterutro enim incurrente%
\indexS{corpus incurrens}
dabitur ipsi \textit{A}, 4
celeritas 1 ad easdem partes,%
\indexS{celeritas post ictum}
licet ipsa corpora \textit{B} et \textit{C}
in ipso \textit{A} idem efficientia,%
\indexS{corpora efficientia idem}
\edtext{%
neque potentiam%
\indexS{potentia corporum concurrentium}
\lbrack neque\rbrack\
quantitatem%
}{%
\lemma{neque}%
\Bfootnote{%
\textit{(1)}~potentiam, neque
\textit{(2)}~potentiam \bars{neque \textit{erg. Hrsg.}} quantitatem%
~\textit{L}%
}}
motus%
\indexS{quantitas motus corporum concurrentium}
eandem habeant.
Sane quantitatem motus eandem non habere
patet,
cum quantitas motus%
\indexS{quantitas motus corporum concurrentium}
ipsius \textit{C} sit 2,
ipsius \textit{B} sit 4.%
\pend%
\vspace{5.0em}%
\pstart%
\noindent%
\centering%
*     *     *
\pend%
\newpage%
\selectlanguage{ngerman}
\frenchspacing%
\begin{ledgroup}
\footnotesize
\pstart
\noindent%
\textls{Ausführliche Datierungsgründe:}%
\pagelabel{41121_dat}
Im Textträger der titellosen und undatierten Aufzeichnung N.\;\ref{41121} ist das Fragment eines Wasserzeichens anzutreffen, das im Leibniz-Nachlass nach heutigem Kenntnisstand lediglich für den Zeitraum 1687 bis 1696 belegt ist. %
Das von der Papiermühle zu Relliehausen am Solling\indexO{Relliehausen am Solling} in Niedersachsen\indexO{Niedersachsen} stammende Papier findet sich mit besonders hohen Häufigkeit in auf die frühen 1690er Jahre datierbaren Handschriften.
Da Leibniz nach seiner Italienreise\indexO{Italien} erst während der ersten Hälfte Juni 1690 in Hannover\indexO{Hannover} eintraf (\cite{01236}\textit{Chronik}, S.\;103), ist es unwahrscheinlich, dass die Aufzeichnung N.\;\ref{41121} vor der zweiten Monatshälfte verfasst wurde.
Als Terminus ante quem der Datierung ist -- unter Vorbehalt -- das Ende der Belegzeit des Wasserzeichens anzunehmen, d.\,h. das Jahr 1696.
Eine spätere Entstehung ist nicht auszuschließen.
\pend
\pstart
Für eine eher späte Datierung der Aufzeichnung N.\;\ref{41121} -- jedenfalls nach Leibnizens Rückkehr nach Hannover\indexO{Hannover} Mitte Juni 1690 -- spricht auch der feststellbare inhaltliche Zusammenhang mit dem wohl auf die Jahre 1699 bis 1701 datierbaren \textit{Essay de dynamique} (\textit{LMG}~VI, S.\;215\textendash231).\cite{02167}\cite{01043} %
Die Gleichungen (1) bis (3) im Abschnitt "Generaliter" von N.\;\ref{41121} (\refpassage{LH_35_09_07_022-023_generaliter-1}{LH_35_09_07_022-023_generaliter-2}) sind nämlich identisch, einschließlich der Wahl der Variabeln, mit denen im \textit{Essay de dynamique} dargelegten (\textit{LMG}~VI, S.\;226\textendash228).\cite{02167}\cite{01043}
Dort werden sie als \textit{trois equations fort simples et belles, et qui contiennent tout ce qui regarde le concours central de deux corps sur une même droite} ausgewiesen (\textit{LMG}~VI, S.\;226).\cite{02167}\cite{01043} 
\pend
\end{ledgroup}
\selectlanguage{latin}%
\frenchspacing%
\normalsize%
